\color{uniblue}{\section{Транспортирование, хранение и консервация}}
\color{black}

Условия транспортировки, хранения (в том числе и в упаковке) и консервации устройств до ввода в эксплуатацию  должны  соответствовать  указанным в таблице \ref{trans:tbl1}. Если требуемые  условия транспортирования и (или) хранения отличаются от приведенных в таблице \ref{trans:tbl1}, то устройство поставляют для условий и сроков, устанавливаемых по ГОСТ~23216 и указываемых в договоре на поставку или заказе-наряде.

\begin{table}[htbp]
\caption{Условия транспортирования и хранения устройств ЮНИТ-М300\hfill\vspace{-0.5\baselineskip}}\label{trans:tbl1}
\centering

\begin{tblr}{
 width=1\textwidth,
 colspec = { X[4,m] X[2,m] X[4,m] X[2,m] X[3.3,m] },
 hlines = {black},
 vlines = {black},
 cell{1}{1} = {r=2,c=1}{c,gray!20},
 cell{1}{2} = {r=1,c=2}{c,gray!20},
 cell{1}{4} = {r=2,c=1}{c,gray!20},
 cell{1}{5} = {r=2,c=1}{c,gray!20},
 cell{2}{2} = {c, gray!20},
 cell{2}{3} = {c, gray!20},
}
 Назначение НКУ & \begin{tabular}{@{}c@{}}Обозначение условий транспортирования \\ в части воздействия\end{tabular} &  & Обозначение условий хранения по ГОСТ~15150 & Допустимые сроки хранения в упаковке и консервации изготовителя, годы \\
  & механических факторов по ГОСТ~23216 & климатических факторов,таких, как условия хранения по ГОСТ~15150 &  & \\
 Для нужд народного хозяйства (кроме районов Крайнего Севера и труднодоступных районов по ГОСТ~15846) & \centering С & \centering 5(ОЖ4) & \centering 1(Л) & \centering 2 \\
 Для нужд народного хозяйства в районы Крайнего Севера и труднодоступные районы по ГОСТ~15846 & \centering С & \centering 5(ОЖ4) & \centering 2(С) & \centering 2 \\
 Для экспорта в макроклиматические районы с умеренным климатом; с тропическим климатом & \centering С & \centering 5(ОЖ4) & \centering 1(Л) & \centering 3 \\
 Для экспорта в макроклиматические районы с тропическим климатом & \centering С & \centering 6(ОЖ2) & \centering 3(Ж3) & \centering 3 \\
\end{tblr}
\end{table}

\vspace{3mm}
\begin{table}[h]
\raggedright
\begin{tabular}{p{10cm}p{5cm}}
Верхнее значение относительной влажности & не более 98\% при плюс 25°С.\\ 
Температура окружающей среды при хранении & от минус 50°С до плюс 50°С. \\
Температура окружающей среды при транспортировании & от минус 60°С до плюс 50°С.\\
\end{tabular}
\end{table}
Условия хранения ЮНИТ-М300 в упаковке у потребителя должны соответствовать условиям хранения 1(Л) по ГОСТ~15150-69.

Расположение ЮНИТ-М300 в хранилищах должно обеспечивать их свободное перемещение и доступ к ним. Необходимо обеспечить отсутствие в окружающей среде складских помещений агрессивных газов в концентрациях, разрушающих металл и изоляцию. 
Расстояние между стенами, полом хранилища и ЮНИТ-М300 должно быть не менее 0,1 м. 
Расстояние  между  отопительными  устройствами  хранилищ  и  ЮНИТ-М300  должно  быть  не  менее 0,5 м. 
Для условий транспортирования в части воздействия механических факторов «С» для экспортных поставок  в  районы  с  умеренным  климатом,  при  наличии  указания  в  заказ-наряде,  допускается 
транспортирование морским путем. 
Требования  по  условиям  хранения  распространяются  на  склады  изготовителя  и  потребителя продукции.

Транспортирование  упакованных  устройств  может  производиться  согласно  ГОСТ~23216 любым видом  закрытого  транспорта,  предохраняющим  изделия  от  воздействия  солнечной  радиации,  резких 
скачков  температур,  атмосферных  осадков  и  пыли  с  соблюдением  мер  предосторожности  против механических воздействий. 
Допустимое число перегрузок упакованных устройств не должно быть более четырех. 
Погрузка,  крепление  и  перевозка  устройств  в  транспортных  средствах  должна  осуществляться  в соответствии  с  действующими  правилами  перевозок  грузов  на  соответствующих  видах  транспорта, причем погрузка, крепление и перевозка устройств железнодорожным транспортом должна производиться в  соответствии  с  «Техническими  условиями  погрузки  и  крепления  грузов»  и  «Правилами  перевозок грузов», утвержденными Министерством путей  сообщения. Все операции с транспортировкой устройств ЮНИТ-М300  необходимо  выполнять  в  соответствии  с  манипуляционными  знаками  маркировки  тары, нанесенной на каждое грузовое место. Упакованные устройства должны быть надежно закреплены для предотвращения их свободного перемещения. 